The main text stated that the distribution of displacements, split up into parallel and orthogonal
components along the rod, can be represented by a multivariate normal distribution.
We first describe in more detail how these distributions are obtained.
Starting from the notation for the vertices $\nu_{ijkl}$ of shape $(N_c,N_v,d,N_t)$ introduced in
Subsection~\ref{subsec:direct-metric}, we focus on one particular cell $c_i$ at a given time point
$t_l$ and denote the resulting quantity with $\nu_{jk}$, therefore dropping the indices $i,l$.
Consider the case where we have a set of exactly known vertices $\nu_{jk}$ and a set of extracted
vertices $\nu_{jk}'$.
The displacements are given by:
\begin{equation}
    \Delta \nu_{jk} = \nu_{jk} - \nu_{jk}'.
\end{equation}
To calculate their positions with respect to the direction of the rod, we assign a
direction $d_{jk}$ at each vertex of the known exact values.
At the tips of the rod this is done by assigning the direction of the closest edge while vertices in
the middle are calculated using a central difference approach:
\begin{align}
    d_{0k} &= \frac{(\nu_{1k}-\nu_{0k})}{||\nu_1 - \nu_0||}\\
    d_{N_v\,k} &= \frac{(\nu_{N_v\,k}-\nu_{N_v-1\,k})}{||\nu_{N_v} - \nu_{N_v-1}||}\\
    d_{jk} &= \frac{(\nu_{j+1\,k}-\nu_{j-1\,k})}{||\nu_{j+1} - \nu_{j-1}||}.
\end{align}
We can use these directions at each vertex to assign rotation matrices
\begin{equation}
    M_{jl}^k = \begin{pmatrix}
        d_{j0} & d_{j1}\\
        -d_{j1} & d_{j0}
    \end{pmatrix}
\end{equation}
at each vertex and apply them to the displacements via
\begin{equation}
    (\Delta\nu_{jk})_\text{directed} = \sum\limits_l M_{jl}^k \Delta\nu_{jl}.
\end{equation}
The resulting points are displayed in Figure~\ref{fig:supplement-extraction-algorithm} (A) which is
identical to the one in the main text.
The main text claims that this distribution is well described by a multivariate normal distribution
with diagonal covariance matrix.
Figure~\ref{fig:supplement-extraction-algorithm} (C,D) show slices of the full distribution along
the parallel and orthogonal axis.
Darker colors indicate earlier stages of the simulation while lighter colors indicate later stages.
Visually, we can already see that the distribution is quite narrow initially but widens with
increased simulation time.
To quantify this behavior, we fit a multivariate normal distribution at each time point to the
distribution of displacements.
The main text has already discussed results for the variances $\sigma_{00}$ and $\sigma_{11}$.
Figure~\ref{fig:supplement-extraction-algorithm} (B) shows the trajectory of the covariance
$\sigma_{01}=\sigma_{10}$ which has been normalized to the value of $\sigma_{00}$.
Its small relative value justifies the statement in the main text and shows no correlation between
parallel and orthogonal displacement values.

Another detail is the resolution of the image used.
Situations of high resolution generally present no concerns but images with low resolutions where a
single bacterium is only made up of a few pixels in its diameter can present problems.
Firstly, the skeletonization algorithm will not be able to accurately determine the middle of the
rod due to the discreteness of the pixel-grid.
Secondly, the same discreteness presents a minimal error threshold which should probably always be
considered.
Within our performed benchmarks, we did not run into such situations since we chose large enough
resolutions.
In general, we do not anticipate that these scenarios are common for most practical applications.

\begin{figure}[H]
    \centering
    \includegraphics[width=\textwidth]{figures/docs/fitting-methods/displacement-calculations-2.pdf}
    \caption{
        (A) Distribution of displacements between exact and extracted vertices parallel and
        orthogonal to the rod.
        (B) Covariance $\sigma_{01}$ normalized by $\sigma_{00}$ over the course of the simulation.
        (C,D) Distributions of displacements sliced parallel (C) and orthogonal (D) to the rod.
        Darker colors indicate earlier stages of the simulation while brighter colors indicate later
        stages.
    }
    \label{fig:supplement-extraction-algorithm}
\end{figure}
